\section{PROBLEM FORMULATION}\label{sec:problem}

\subsection{Setting}
A pipeline emits a telemetry snapshot $s_k$ at every control tick $k$, with tick period
$\tau = 15$\,s. A fault $f$ is injected at time $t_f$, detected at $t_d \ge t_f$, and resolved at
$t_r \ge t_d$, either through an executed mitigation or an escalation the on-call engineer handles. A
\emph{configuration} $\pi$ is a responder mapping snapshots to actions. For an agent configuration,
each enabled agent $g$ proposes an action
$a = (g, \text{type}, \text{target}, \text{params}, \text{justification}, \text{confidence})$ drawn
from the typed action space $\mathcal{A}$ of the contract (Section~\ref{sec:system}). A policy gate
$G$ decides each proposal from a context $c = (m, B, n, \kappa, v)$ that the system itself builds,
beyond the agent's influence: the action's projected marginal cost $m$, the remaining budget $B$,
the number $n$ of actions taken in the last ten minutes, the schema-compatibility class $\kappa$,
and a flag $v$ marking whether a prior version exists. The gate returns one of \textsc{allow},
\textsc{deny}, \textsc{escalate}, or \textsc{allow-and-notify}.

\subsection{Gate Semantics}
The gate checks the rate limit of Eq.~\eqref{eq:rate} before any agent-specific policy and defaults
toward the human on failure:
\begin{equation}
  G(a,c) = \textsc{deny} \quad \text{if } n \ge n_{\max}, \qquad n_{\max}=5,
  \label{eq:rate}
\end{equation}
and any proposal that matches no policy pack, or that cannot be evaluated because the policy engine
is unreachable, is escalated rather than allowed. An optimization action that scales resources
clears the cost policy only if the condition in Eq.~\eqref{eq:budget} holds:
\begin{equation}
  m \le 0 \quad \text{or} \quad m \le B,
  \label{eq:budget}
\end{equation}
so a scale-down, carrying non-positive marginal cost, always clears. Rollback is permitted only when
$v$ holds, and a schema change is permitted only when $\kappa$ is backward-compatible.
Section~\ref{sec:system} describes all four policy packs.

\subsection{Metrics}
Each run injects a single fault. Its recovery time, cost, and decision correctness are given by Eqs.~\eqref{eq:mttr},
\eqref{eq:cost}, and~\eqref{eq:decision}:
\begin{align}
  \mathrm{MTTR}(r) &= t_r - t_f, \label{eq:mttr}\\
  C(r) &= C_{\mathrm{meas}}(r) + \rho_c\, u_\pi\, T_h, \label{eq:cost}\\
  D(r) &= \mathbf{1}\bigl[\text{executed mitigation} \in \mathcal{A}^{*}(f)\bigr], \label{eq:decision}
\end{align}
where $C_{\mathrm{meas}}$ is the metered compute and storage cost, with compute-unit-seconds priced
at $\rho_c = 0.05$ per unit-second and storage at 0.01 per GB-hour; $u_\pi$ is the allocation
configuration $\pi$ holds over the horizon $T_h = 300$\,s; and $\mathcal{A}^{*}(f)$ is the accepted
mitigation set for the injected fault. The three configurations that right-size,
\texttt{autoscale}, \texttt{optimization\_only}, and \texttt{full}, hold $u_\pi = 3$ units, and the
rest hold a static $u_\pi = 8$. Non-agent configurations never execute an agent mitigation, so
$D(r)$ in Eq.~\eqref{eq:decision} is $0$ for them by construction. Our MTTR in Eq.~\eqref{eq:mttr} includes detection lag, since $t_f$ here is the
injection time, whereas the base paper measures recovery from detection to resumption. Manual
interventions and data freshness are direct counts and lags read from telemetry.

\subsection{The Evaluation-Validity Question}
Let $\tilde{x}_{\mathrm{live}}$ and $\tilde{x}_{\mathrm{mock}}$ denote the median of metric $x$ for
the same configuration under a real language model and under a deterministic mock, respectively. We
define the divergence ratio as
\begin{equation}
  R_x = \frac{\tilde{x}_{\mathrm{live}}}{\tilde{x}_{\mathrm{mock}}}.
  \label{eq:ratio}
\end{equation}
Reporting $\tilde{x}_{\mathrm{mock}}$ as the system's performance is valid only when $R_x$ of
Eq.~\eqref{eq:ratio} sits close to one. Section~\ref{sec:livemock} measures $R_x$ for every metric, and RQ4 of
Section~\ref{sec:intro} asks how far it strays from one.

\subsection{Threat Model}
The adversary controls only its own proposal $a$. The context $c$, the policy engine, and the
executor are trusted, system-built components. The safety property under test is containment: no
proposal that the documented policy semantics classifies as unsafe may be allowed.

\begin{algorithm}[t]
\caption{One control tick.}
\label{alg:tick}
\begin{algorithmic}[1]
\STATE record a heartbeat
\IF{the kill switch is engaged} \RETURN \ENDIF
\STATE run the monitoring agent: detect faults and stamp $t_d$
\IF{an unresolved fault exists}
  \STATE $P \leftarrow \emptyset$
  \FOR{each enabled reactive agent $g \in \{\text{schema}, \text{recovery}, \text{optimization}\}$}
    \STATE $a \leftarrow$ observe and reason (invalid output becomes \texttt{no\_action})
    \STATE add $a$ to $P$ unless it is \texttt{no\_action}
  \ENDFOR
  \STATE for each contested target, keep the proposal with the highest (agent priority, confidence)
  \FOR{each winning proposal $a$}
    \STATE skip if the target lock is held or the blast-radius cap is reached
    \STATE evaluate $d \leftarrow G(a,c)$, then commit an intent row
    \STATE execute $a$ if $d$ allows it, or hand it to the on-call if $d$ escalates
    \STATE update the audit row with the outcome
  \ENDFOR
\ENDIF
\end{algorithmic}
\end{algorithm}

Algorithm~\ref{alg:tick} summarizes the control loop Section~\ref{sec:system} describes in prose.
Agent priority orders recovery above schema above optimization, and the non-agent configurations
skip the loop entirely, resolving faults directly instead (Section~\ref{sec:method}).
